\documentclass[11pt]{article}
\usepackage[T1]{fontenc}
\usepackage[utf8]{inputenc}
\usepackage{lmodern}
\usepackage[a4paper,margin=27mm]{geometry}
\usepackage{amsmath,amssymb,amsthm,booktabs,longtable,listings,microtype}
\usepackage[hidelinks]{hyperref}
\hypersetup{pdftitle={All partition-restricted counts of nondegenerate 2 by 4 by 3 hypermatrices},
  pdfauthor={Xiaojian Zeng}}
\newtheorem{theorem}{Theorem}
\newtheorem{lemma}{Lemma}
\newcommand{\F}{\mathbb F}
\newcommand{\Z}{\mathbb Z}
\newcommand{\inv}{\operatorname{inv}}
\newcommand{\GL}{\operatorname{GL}}
\newcommand{\rank}{\operatorname{rank}}
\newcommand{\qint}[1]{[#1]_q}



\title{All partition-restricted counts of nondegenerate\\$2\times4\times3$ hypermatrices}
\author{Xiaojian Zeng\\[3pt]
  \small PBC School of Finance, Tsinghua University, 100190, China\\
  \small\texttt{xiaojianzeng@foxmail.com}}
\date{}
\begin{document}
\maketitle
\begin{abstract}
We prove the format $2\times4\times3$ specialization of the
Koprowski--Lewis product conjecture for nondegenerate hypermatrices over
finite fields with partition-shaped zero restrictions. The proof evaluates
their augmented-cell sum for all 55 admissible pairs of partitions and
all prime powers. An explicit order for coefficient-one elimination makes
the cell parametrizations valid in every characteristic. The finite part
consists of 55 integer polynomial identities, for which we give the complete
coefficient table and an executable census. Separate sparse and dense
polynomial computations verify the elimination certificates; exhaustive
binary enumeration supplies an additional check. The result includes the
40 shapes outside the two previously proved product families in this
format. It makes no assertion for larger formats.
\end{abstract}
\noindent\textbf{Keywords:} hyperdeterminant; finite field; restricted enumeration;
partition; exact computation.

\section{Statement and relation to prior work}

Let $q$ be a prime power and let $\lambda,\mu$ be weakly decreasing triples
of nonnegative integers satisfying
\begin{equation}\label{eq:shapes}
 \mu_j\leq\lambda_j,\qquad \lambda_j\leq4-j,\qquad
 \mu_j\leq3-j\qquad(1\leq j\leq3).
\end{equation}
For $P=(\lambda,\mu)$, an ordered pair of matrices
$A,B\in\F_q^{4\times3}$ \emph{respects $P$} if
\begin{equation}\label{eq:zeros}
 A_{rj}=0\quad\text{when }r>4-\lambda_j,\qquad
 B_{rj}=0\quad\text{when }r>4-\mu_j.
\end{equation}
Every unrestricted entry may take any value, including zero. We count
arrays individually, without identifying them under a group action.
Nondegeneracy means that the boundary-format hyperdeterminant is nonzero.
In this format it is equivalent to
\begin{equation}\label{eq:rank}
 \rank(sA+tB)=3\quad\text{for every }(s,t)\in
 \overline{\F}_q^{\,2}\setminus\{(0,0)\};
\end{equation}
see Aitken's rank criterion, as stated in
\cite[Lemma 2.1]{KL} (Lemma 2.5 in the arXiv version).
The algebraic closure in \eqref{eq:rank} is essential.
Write $N_P(q)$ for this count and $[r]_q=\sum_{a=0}^{r-1}q^a$.

\begin{theorem}\label{thm:main}
For every prime power $q$ and every pair satisfying \eqref{eq:shapes},
\begin{equation}\label{eq:main}
 N_P(q)=q^9(q-1)^6
 \prod_{j=1}^3[5-j-\lambda_j]_q
 \prod_{j=1}^3[4-j-\mu_j]_q.
\end{equation}
\end{theorem}

Haglund's interpretation of $q$-rook polynomials counts matrices of
prescribed rank supported on Ferrers boards
\cite[Theorem 1 in the preprint]{Haglund}.
Aitken obtained the unrestricted boundary-format hypermatrix count from
a free transitive action modulo scalar matrices \cite{Aitken}.
Koprowski and Lewis refined that enumeration to the restrictions
\eqref{eq:zeros} and proposed the corresponding product for arbitrary
$2\times(k+1)\times k$ format \cite[Conjecture 3.1]{KL}.
Their Theorems 3.2 and 3.3 establish the cases $\mu=0$ and the maximal
shape; their Section 3.3 also treats every shape for $k=2$.
All numbering here refers to the journal version of \cite{KL}, unless
otherwise specified.

For $k=3$, the $\mu=0$ family has 14 shapes, and the maximal shape
$((3,2,1),(2,1,0))$ is a further case. Section 3.3 of \cite{KL}
identifies 40 other shapes without giving their complete product
evaluation. Theorem~\ref{thm:main} supplies that evaluation, using exactly
the permutation sum proposed in their Section 5.1. The orbit reduction,
the pure-power cell count, and polynomial positivity are established
ingredients. Our contribution is the complete finite calculation in this
format, with explicit integral elimination and reproducible certificates.
The comparison is with the inspected journal and arXiv versions and a
targeted literature search through 12 September 2026; it is not a claim of
exhaustive priority verification. The general-$k$ conjecture is outside
the theorem.

\section{The augmented-cell reduction}

Let $e_1,\ldots,e_4$ be the standard basis and put
$E_0=(e_1,e_2,e_3)$ and $E_1=(e_2,e_3,e_4)$.
For $\sigma\in S_4$ and $\pi\in S_3$, introduce distinct variables
$x_{ab}$ for inversions $a<b$, $\sigma(a)>\sigma(b)$, and $y_{ij}$
for inversions $i<j$, $\pi(i)>\pi(j)$. Define $L$ and $R$ by
\begin{align}
 L_{\sigma(a),a}&=1,&
 L_{\sigma(a),b}&=x_{ab}
       &&(a<b,\ \sigma(a)>\sigma(b)),\label{eq:L}\\
 R_{\pi(i),i}&=1,&
 R_{\pi(j),i}&=y_{ij}
       &&(i<j,\ \pi(i)>\pi(j)),\label{eq:R}
\end{align}
with every other entry zero. Thus $L$ is $4\times4$, $R$ is $3\times3$,
and the two augmented slices are
\begin{equation}\label{eq:H}
 H_0=LE_0R,\qquad H_1=LE_1R.
\end{equation}
These are the conventions of \cite[Definitions 4.5--4.7]{KL}.
In the group-action notation of that paper the second acting matrix is
$R^{\mathsf T}$; $R$ itself multiplies on the right in \eqref{eq:H}.
There are $n=\inv(\sigma)+\inv(\pi)$ variables.

Let $C_{\sigma,\pi}(P,q)$ count their assignments in $\F_q$ for which
\eqref{eq:H} respects $P$. The established reduction is
\begin{equation}\label{eq:reduction}
 N_P(q)=q^9(q-1)^6
       \sum_{\sigma\in S_4,\,\pi\in S_3}C_{\sigma,\pi}(P,q).
\end{equation}
This is \cite[Corollary 4.2]{KL}. We recall the normalization to fix the
counting convention. Aitken's action of
\[
 (\GL_4(\F_q)\times\GL_3(\F_q))/
 \{(cI_4,c^{-1}I_3):c\in\F_q^\times\}
\]
on nondegenerate arrays is free and transitive \cite{Aitken}.
Bruhat decomposition, in the form used in \cite[Theorem 4.3]{KL},
separates the augmented matrices from two triangular factors of orders
$q^6(q-1)^4$ and $q^3(q-1)^3$. Division by the $q-1$ scalar
representatives gives $q^9(q-1)^6$.
Left upper triangular multiplication preserves each column's forbidden
bottom segment. Right upper triangular multiplication preserves the zero
pattern because permitted column heights are nondecreasing. The inverse
operations have the same properties. Finally, distinct normalized
assignments within one augmented cell give distinct arrays: equality
would imply a scalar relation between the acting matrices, and the
prescribed unit pivots force this scalar to be one. These observations
explain why \eqref{eq:reduction} counts ordered arrays with the stated
restrictions.

\section{Integral elimination in each cell}

The pure-power conclusion below is the counting principle of
\cite[Theorem 4.4 and Section 5.1]{KL}. We give its explicit elimination
order here to explain why the finite certificate applies in every
characteristic.

\begin{lemma}\label{lem:cell}
The augmented cell is empty if one of the six positions
\begin{equation}\label{eq:pivots}
 (\sigma(\pi(c)),c)\text{ in }H_0,\qquad
 (\sigma(\pi(c)+1),c)\text{ in }H_1\qquad(1\leq c\leq3)
\end{equation}
is forbidden. Otherwise let $h$ be the number of forbidden entries of
\eqref{eq:H} that are nonzero polynomials over $\Z$. Then
\begin{equation}\label{eq:cell}
 C_{\sigma,\pi}(P,q)=q^{n-h}
\end{equation}
for every prime power $q$.
\end{lemma}

\begin{proof}
Each entry in \eqref{eq:pivots} is exactly 1. Indeed, in the front slice
the two triangular supports meet only at the intermediate index $\pi(c)$;
in the back slice the same argument uses $\pi(c)+1$ on the left.
These are the only entries with a constant term. A forbidden pivot is
therefore impossible to cancel. If every pivot is allowed, setting all
variables to zero respects the shape.

Assume henceforth that the pivots are allowed. For a forbidden front
entry $(r,c)$ put $a=\sigma^{-1}(r)$ and $v=\pi(c)$. Contributing
intermediate indices satisfy $a\leq t\leq v$. If $a>v$ the entry is
zero; $a=v$ would be a forbidden pivot. If $a<v$, the allowed pivot
in column $c$ and the forbidden row $r$ imply $r>\sigma(v)$, so
$x_{av}$ exists. Put $b=\pi^{-1}(a)$. We have $b>c$: otherwise the
front pivot in the earlier column $b$ would give
$r\leq4-\lambda_b\leq4-\lambda_c$, a contradiction.
Consequently the equation for this entry is
\begin{equation}\label{eq:front}
 y_{cb}+x_{av}
   +\sum_{\substack{a<t<v\\\text{both variables exist}}}
          x_{at}y_{c,\pi^{-1}(t)}=0.
\end{equation}
Its target variable is $y_{cb}$, which occurs only in the displayed
linear term.

For a forbidden back entry, again put $a=\sigma^{-1}(r)$ and
$v=\pi(c)$, and set $d=v+1$. If $a>d$ the entry is zero; $a=d$
would be a forbidden pivot. For $a<d$ we have $r>\sigma(d)$, so the
entry contains $x_{ad}$ with coefficient one. Every other term uses
only $x_{a,t+1}$ with $t<v$, a variable $y_{cj}$, or a product of
such variables. We take $x_{ad}$ as its target.

Order the variables lexicographically by
\begin{equation}\label{eq:order}
 K(x_{ab})=(b,0,0,a),\qquad
 K(y_{ij})=(\pi(i),1,-\pi(j),0).
\end{equation}
In \eqref{eq:front}, $x_{av}$ precedes the target at the second key
coordinate, and $x_{at}$ at the first. The intermediate
$y_{c,\pi^{-1}(t)}$ precedes it at the third coordinate because $t>a$.
In a back equation the target has first coordinate $v+1$, whereas
every other variable has first coordinate at most $v$.

The front targets are distinct: $(c,b)$ determines the column and the
row $\sigma(\pi(b))$. The back targets are distinct because $(a,d)$
determines the row and the column $\pi^{-1}(d-1)$. The two target
sets use disjoint variable types. Thus the $h$ nonzero equations have
$h$ distinct targets and each has the form $z+f=0$, where $f$ depends
only on earlier variables. Assign the remaining $n-h$ variables
freely and solve in order \eqref{eq:order}. This is a bijection, with
inverse the projection to the free coordinates. All solved expressions
are integer polynomials; no division or exceptional characteristic
enters the argument. This proves \eqref{eq:cell}.
\end{proof}

\section{The complete finite evaluation}

The following ranges enumerate exactly \eqref{eq:shapes}:
\begin{align*}
 0\leq\lambda_1&\leq3,&
 0\leq\lambda_2&\leq\min(\lambda_1,2),&
 0\leq\lambda_3&\leq\min(\lambda_2,1),\\
 0\leq\mu_1&\leq\min(\lambda_1,2),&
 0\leq\mu_2&\leq\min(\mu_1,\lambda_2,1),&\mu_3&=0.
\end{align*}
They give 55 shapes. For each one we use every pair in $S_4\times S_3$,
with no symmetry quotient. Discard the pairs with a forbidden pivot and
let $c_d(P)$ be the number of remaining pairs with $n-h=d$.
Table~\ref{tab:census} gives every coefficient vector, in ascending
powers of $q$. In each row it is equal to the coefficient vector of
\begin{equation}\label{eq:product}
 \prod_{i=1}^{6}(1+q+\cdots+q^{r_i-1}),\qquad
 (r_1,\ldots,r_6)=
 (4-\lambda_1,3-\lambda_2,2-\lambda_3,
  3-\mu_1,2-\mu_2,1-\mu_3).
\end{equation}
All six lengths are positive. This is coefficient comparison in
$\Z[q]$, rather than evaluation at selected field sizes.

The short program in Appendix~\ref{app:code} reproduces the census.
For this purpose, nonzero support of a generic entry can be detected
by setting every inversion variable to 1 and multiplying over the
ordinary integers. The matrices in \eqref{eq:L}--\eqref{eq:R} have
only positive monomial coefficients, and multiplication introduces
no subtraction. Hence an entry is the zero polynomial if and only if
this integer evaluation is zero. The evaluation is used only to count
$h$; the pivot test and Lemma~\ref{lem:cell} supply the field-independent
solution count. The right side of \eqref{eq:product} is computed by
integer convolution.

The full polynomial certificate separately reconstructs all 3,456
entries in the 144 augmented arrays over $\Z$. Among the 7,920
shape--cell cases, 6,759 have a forbidden unit entry and 1,161 survive.
The surviving cases contain 997 constrained nonzero entries. Their
targets are distinct, the order \eqref{eq:order} is valid, and
substitution of the serialized solutions gives zero in every original
equation. The largest solution-polynomial degree is 3. A second
implementation using dense symbolic matrices in SymPy~1.14.0
\cite{SymPy} reconstructs the entries, checks every elimination
identity, and verifies all 55 product identities without importing the
sparse implementation.

Separating exhaustive coverage from exact certificate checks is useful
in finite computer-assisted proofs. For example, Zeng, Liu, Ratnavelu
and Ong \cite{ZengCertificates} reconstruct the finite
objects and verify integer predicates in their ternary-geometry
computations. Here the corresponding obligations are complete
shape/permutation coverage, integral elimination, and equality of
integer coefficient vectors; the geometry results are not assumptions
of our argument.

\begin{proof}[Proof of Theorem~\ref{thm:main}]
By \eqref{eq:reduction} and Lemma~\ref{lem:cell}, the normalized count
$N_P(q)/(q^9(q-1)^6)$ is $\sum_d c_d(P)q^d$.
The complete census identifies this polynomial with \eqref{eq:product}
for every shape in \eqref{eq:shapes}. Its six factors are precisely the
six $q$-integers in \eqref{eq:main}, proving the assertion for every
prime power.
\end{proof}

As boundary checks, the unrestricted coefficient vector is
$(1,5,13,23,30,30,23,13,5,1)$, representing $[4]!_q[3]!_q$, where
$[r]!_q=\prod_{i=1}^r[i]_q$. The maximal shape has coefficient vector
$(1)$ and exactly one surviving cell. Both cases are included in the
census, not used as substitutes for it.

\section{Binary verification and reproducibility}

A separate exhaustive test covers the $2^{24}$ ordered binary arrays.
The necessary full-rank conditions on $A$ and $B$ reduce this to
$2520^2=6{,}350{,}400$ pairs. For each pair the program computes the
greatest common divisor in $\F_2[t]$ of the four maximal minors of
$A+tB$. A gcd equal to 1 rules out a finite algebraic-closure root
common to all minors; rank 3 of $B$ handles the point at infinity.
Thus the test implements \eqref{eq:rank} exactly. Integer array
operations use NumPy~2.2.6 \cite{NumPy}.

All 55 binary counts agree with \eqref{eq:main}; the unrestricted count
is $3{,}386{,}880$. Testing only the three rational pencil directions
would admit $403{,}200$ additional arrays. For instance,
\[
 A=\begin{pmatrix}0&0&1\\0&1&0\\1&0&0\\0&0&0\end{pmatrix},\qquad
 B=\begin{pmatrix}1&1&0\\0&0&1\\0&1&0\\0&0&0\end{pmatrix}
\]
have rank 3 in all three rational directions, but their only nonzero
maximal minor is $1+t^2+t^3$, which has a root over $\overline{\F}_2$.
This binary test is corroboration; the all-$q$ proof is the integral
argument and the finite polynomial census above.

The companion directory includes both polynomial verifiers, all cell
certificates, and the binary enumeration program with its recorded
counts. The accompanying reproduction guide specifies their filenames,
directory layout and commands. The sparse verifier and the census in
Appendix~\ref{app:code} use only the Python standard library.
The dense reconstruction and the binary enumeration use the software
versions specified above. The dense verification did not rerun the
full binary enumeration; it did check its recorded counts against the
polynomials and independently recompute the diagnostic witness.

\paragraph{Computational and AI assistance.}
OpenAI Codex assisted with literature searches, development of the
elimination argument, writing and executing the verification programs,
and manuscript preparation. A separate model-context review wrote and
ran the dense symbolic reconstruction. This provides a separate
implementation check, not external peer review or proof-assistant
certification. The mathematical conclusions rest on the stated
dependencies, the integral elimination argument, and the inspectable
exact computations.

\appendix
\section{Executable coefficient census}\label{app:code}
This standard-library Python program prints the 55 vectors in
Table~\ref{tab:census}, padded with trailing zeros to length ten.
Permutations store one-based values; matrix indices in the code are
zero-based. Its integer template entries are evaluations at 1, as
justified in Section~4. Running \texttt{python3 manuscript/census.py}
from the companion directory reproduces the output.
\begin{lstlisting}[language=Python,basicstyle=\ttfamily\footnotesize,columns=fullflexible,keepspaces=true,showstringspaces=false,frame=single,breaklines=true,tabsize=4]
"""Exact coefficient census; run with Python 3 (standard library only)."""
from itertools import permutations, product


def inv(p):
    return sum(p[i] > p[j] for i in range(len(p))
               for j in range(i + 1, len(p)))


def multiply(a, b):
    return [[sum(x * y for x, y in zip(row, col))
             for col in zip(*b)] for row in a]


def templates(s, p):
    left = [[0] * 4 for _ in range(4)]
    right = [[0] * 3 for _ in range(3)]
    for a in range(4):
        left[s[a] - 1][a] = 1
        for b in range(a + 1, 4):
            left[s[a] - 1][b] = int(s[a] > s[b])
    for i in range(3):
        right[p[i] - 1][i] = 1
        for j in range(i + 1, 3):
            right[p[j] - 1][i] = int(p[i] > p[j])
    return [multiply([row[f:f + 3] for row in left], right)
            for f in range(2)]


def census():
    pairs = list(product(permutations(range(1, 5)),
                         permutations(range(1, 4))))
    data = [(s, p, inv(s) + inv(p), templates(s, p))
            for s, p in pairs]
    rows = []
    for six in product(range(4), repeat=6):
        lam, mu = six[:3], six[3:]
        if not (lam[0] >= lam[1] >= lam[2]
                and mu[0] >= mu[1] >= mu[2]):
            continue
        if not all(mu[j] <= lam[j] <= 3 - j
                   and mu[j] <= 2 - j for j in range(3)):
            continue
        counts = [0] * 10
        for s, p, n, hmat in data:
            if any(s[p[c] - 1 + f] > 4 - sh[c]
                   for f, sh in enumerate((lam, mu))
                   for c in range(3)):
                continue
            h = sum(hmat[f][r][c] != 0
                    for f, sh in enumerate((lam, mu))
                    for r in range(4) for c in range(3)
                    if r + 1 > 4 - sh[c])
            assert 0 <= n - h <= 9
            counts[n - h] += 1
        lengths = [4 - j - lam[j] for j in range(3)]
        lengths += [3 - j - mu[j] for j in range(3)]
        coeff = [1]
        for length in lengths:
            out = [0] * (len(coeff) + length - 1)
            for i, a in enumerate(coeff):
                for j in range(length):
                    out[i + j] += a
            coeff = out
        assert counts == coeff + [0] * (10 - len(coeff))
        rows.append((lam, mu, counts))
    assert len(rows) == 55 and len(pairs) == 144
    return rows


if __name__ == "__main__":
    for lam, mu, counts in census():
        print(lam, mu, counts)
\end{lstlisting}

\clearpage
\section{All coefficient vectors}
In Table~\ref{tab:census}, a string such as $210$ denotes the triple
$(2,1,0)$. The column $s$ is the number of surviving permutation pairs,
and equals the sum of the displayed coefficients. Trailing zero
coefficients are omitted. The factor lengths are determined by
\eqref{eq:product}.
{
\small

\begin{longtable}{@{}ccrl@{}}
\caption{Complete coefficient census for $N_P(q)/(q^9(q-1)^6)$.}\label{tab:census}\\
\toprule
$\lambda$ & $\mu$ & $s$ & $(c_0,c_1,\ldots)$ \\
\midrule
\endfirsthead
\multicolumn{4}{c}{Table \ref{tab:census}\ (continued)}\\
\toprule
$\lambda$ & $\mu$ & $s$ & $(c_0,c_1,\ldots)$ \\
\midrule
\endhead
\midrule
\multicolumn{4}{r}{Continued on next page}\\
\endfoot
\bottomrule
\endlastfoot
$000$ & $000$ & 144 & $(1,5,13,23,30,30,23,13,5,1)$ \\
$100$ & $000$ & 108 & $(1,5,13,22,26,22,13,5,1)$ \\
$100$ & $100$ & 72 & $(1,5,12,18,18,12,5,1)$ \\
$110$ & $000$ & 72 & $(1,5,12,18,18,12,5,1)$ \\
$110$ & $100$ & 48 & $(1,5,11,14,11,5,1)$ \\
$110$ & $110$ & 24 & $(1,4,7,7,4,1)$ \\
$111$ & $000$ & 36 & $(1,4,8,10,8,4,1)$ \\
$111$ & $100$ & 24 & $(1,4,7,7,4,1)$ \\
$111$ & $110$ & 12 & $(1,3,4,3,1)$ \\
$200$ & $000$ & 72 & $(1,5,12,18,18,12,5,1)$ \\
$200$ & $100$ & 48 & $(1,5,11,14,11,5,1)$ \\
$200$ & $200$ & 24 & $(1,4,7,7,4,1)$ \\
$210$ & $000$ & 48 & $(1,5,11,14,11,5,1)$ \\
$210$ & $100$ & 32 & $(1,5,10,10,5,1)$ \\
$210$ & $110$ & 16 & $(1,4,6,4,1)$ \\
$210$ & $200$ & 16 & $(1,4,6,4,1)$ \\
$210$ & $210$ & 8 & $(1,3,3,1)$ \\
$211$ & $000$ & 24 & $(1,4,7,7,4,1)$ \\
$211$ & $100$ & 16 & $(1,4,6,4,1)$ \\
$211$ & $110$ & 8 & $(1,3,3,1)$ \\
$211$ & $200$ & 8 & $(1,3,3,1)$ \\
$211$ & $210$ & 4 & $(1,2,1)$ \\
$220$ & $000$ & 24 & $(1,4,7,7,4,1)$ \\
$220$ & $100$ & 16 & $(1,4,6,4,1)$ \\
$220$ & $110$ & 8 & $(1,3,3,1)$ \\
$220$ & $200$ & 8 & $(1,3,3,1)$ \\
$220$ & $210$ & 4 & $(1,2,1)$ \\
$221$ & $000$ & 12 & $(1,3,4,3,1)$ \\
$221$ & $100$ & 8 & $(1,3,3,1)$ \\
$221$ & $110$ & 4 & $(1,2,1)$ \\
$221$ & $200$ & 4 & $(1,2,1)$ \\
$221$ & $210$ & 2 & $(1,1)$ \\
$300$ & $000$ & 36 & $(1,4,8,10,8,4,1)$ \\
$300$ & $100$ & 24 & $(1,4,7,7,4,1)$ \\
$300$ & $200$ & 12 & $(1,3,4,3,1)$ \\
$310$ & $000$ & 24 & $(1,4,7,7,4,1)$ \\
$310$ & $100$ & 16 & $(1,4,6,4,1)$ \\
$310$ & $110$ & 8 & $(1,3,3,1)$ \\
$310$ & $200$ & 8 & $(1,3,3,1)$ \\
$310$ & $210$ & 4 & $(1,2,1)$ \\
$311$ & $000$ & 12 & $(1,3,4,3,1)$ \\
$311$ & $100$ & 8 & $(1,3,3,1)$ \\
$311$ & $110$ & 4 & $(1,2,1)$ \\
$311$ & $200$ & 4 & $(1,2,1)$ \\
$311$ & $210$ & 2 & $(1,1)$ \\
$320$ & $000$ & 12 & $(1,3,4,3,1)$ \\
$320$ & $100$ & 8 & $(1,3,3,1)$ \\
$320$ & $110$ & 4 & $(1,2,1)$ \\
$320$ & $200$ & 4 & $(1,2,1)$ \\
$320$ & $210$ & 2 & $(1,1)$ \\
$321$ & $000$ & 6 & $(1,2,2,1)$ \\
$321$ & $100$ & 4 & $(1,2,1)$ \\
$321$ & $110$ & 2 & $(1,1)$ \\
$321$ & $200$ & 2 & $(1,1)$ \\
$321$ & $210$ & 1 & $(1)$ \\
\end{longtable}
}

\small
\bibliographystyle{plain}
\bibliography{references}
\end{document}
